L'IMU étant en coordonnées cartésienne et le récepteur GNSS en coordonnées sphériques, il faut les harmoniser avant de les fusionner

Nous avons pu remarqué lors de notre phase d'analyse, que l'absence de cette harmonisation pouvait affecter la représentation finale

Nous avons donc implémenté deux fonctions de conversion:
\[
\left\{
\begin{array}{l}
    x = (R + Altitude) * \sin(\theta) * \cos(\phi) \\
    y = (R + Altitude) * \sin(\theta) * \sin(\phi) \\
    y = (R + Altitude) * \cos(\theta)
\end{array}
\right.
\]

\[
\left\{
\begin{array}{l}
    \theta = \arccos(\frac{z}{\sqrt{x^2 + y^2 +z^2}}) \\
    \phi = \arctan(\frac{y}{x}) \\
    Altitude = \sqrt{x^2 + y^2 +z^2} - R
\end{array}
\right.
\]

Ainsi, grâce au multiprocessing nous avons des acquisition qui arrivent au fur et à mesures ce qui nous permet de Synchroniser les acquisitions,
et des fonctions de conversion afin d'avoir un filtrage "cohérent" par rapport aux équations dynamiques.